\documentclass[a4paper,10pt]{extarticle}
\usepackage[utf8]{inputenc}
\usepackage[spanish]{babel}
\usepackage[margin=0.5cm]{geometry}
\usepackage{multicol}
\usepackage{amsmath}
\usepackage{enumitem}
\usepackage{titlesec}
\usepackage{mathptmx}
\usepackage{graphicx}
\usepackage{csquotes}

% Configuración de listas y títulos
\setlist{nosep}
\titlespacing{\section}{0pt}{1pt}{1pt}
\setlist[enumerate]{leftmargin=*, topsep=1pt, itemsep=2pt}
\setlength{\parindent}{0pt}

\begin{document}

% Encabezado
\begin{center}
    {\large \textbf{SIMULACRO DE EXAMEN: TELEDETECCIÓN Y PROCESADO DE IMÁGENES}}\\
\end{center}

\begin{multicols*}{2}

% ---------------------------
% PARTE I: TEST
% ---------------------------

\section*{I. Teoría (4 ptos) - 0.4 c/u}
\textit{\footnotesize Seleccione la respuesta única correcta.}

\begin{enumerate}

    % TEMA 1: Modalidades / Introducción
    \item \textbf{¿Cuál de las siguientes modalidades de imagen médica utiliza un radiofármaco inyectado en el paciente para obtener información funcional o metabólica?}
    \begin{enumerate}[label=\alph*)]
        \item Resonancia Magnética (RM).
        \item Tomografía Computarizada (TC).
        \item Tomografía por Emisión de Positrones (PET).
        \item Radiografía Simple.
    \end{enumerate}

    % TEMA 1: Calidad de imagen
    \item \textbf{El \textquote{Efecto de Volumen Parcial} ocurre cuando:}
    \begin{enumerate}[label=\alph*)]
        \item El paciente se mueve durante la adquisición.
        \item Un vóxel contiene tejidos diferentes y su intensidad es un promedio de ellos.
        \item Hay objetos metálicos dentro del campo de visión.
        \item La relación señal a ruido es demasiado baja.
    \end{enumerate}

    % TEMA 2: Histograma / Transformaciones
    \item \textbf{¿Qué transformación de intensidad utilizaría para redistribuir los niveles de gris de una imagen oscura y aprovechar todo el rango dinámico disponible?}
    \begin{enumerate}[label=\alph*)]
        \item Umbralización binaria.
        \item Inversión de color (negativo).
        \item Ecualización del histograma.
        \item Filtro de mediana.
    \end{enumerate}

    % TEMA 2: Filtros / Espacio-Frecuencia
    \item \textbf{El filtro de \textquote{Difusión Anisotrópica} se caracteriza principalmente por:}
    \begin{enumerate}[label=\alph*)]
        \item Suavizar la imagen preservando los bordes.
        \item Eliminar únicamente el ruido impulsivo ("sal y pimienta").
        \item Aumentar el ruido para resaltar texturas.
        \item Ser un filtro lineal equivalente al filtro de media.
    \end{enumerate}

    % TEMA 3: Transformaciones Geométricas
    \item \textbf{En una transformación geométrica, ¿qué operación modifica la forma de la imagen convirtiendo un cuadrado en un paralelogramo (sesgo)?}
    \begin{enumerate}[label=\alph*)]
        \item Traslación.
        \item Cizallamiento (Shear).
        \item Rotación.
        \item Escalado.
    \end{enumerate}

    % TEMA 3: Registro
    \item \textbf{¿Qué diferencia principal existe entre un registro rígido y un registro afín?}
    \begin{enumerate}[label=\alph*)]
        \item El registro afín permite deformaciones locales elásticas.
        \item El registro rígido solo permite rotación y traslación; el afín añade escalado y cizallamiento.
        \item El registro rígido tiene más grados de libertad que el afín.
        \item No hay diferencia, son sinónimos.
    \end{enumerate}

    % TEMA 4: Segmentación (Umbralización)
    \item \textbf{El método de Otsu para la umbralización automática busca el umbral óptimo que:}
    \begin{enumerate}[label=\alph*)]
        \item Minimiza la varianza inter-clase.
        \item Maximiza la varianza inter-clase (o minimiza la intra-clase).
        \item Selecciona siempre el valor medio del histograma (ej. 128).
        \item Elimina los píxeles conectados al borde de la imagen.
    \end{enumerate}

    % TEMA 4: Segmentación (Contornos Activos)
    \item \textbf{En los modelos de Contornos Activos (Snakes), la \textquote{energía interna} de la curva se encarga de:}
    \begin{enumerate}[label=\alph*)]
        \item Atraer la curva hacia los bordes de la imagen.
        \item Controlar la suavidad y elasticidad de la curva (continuidad).
        \item Calcular el gradiente de la imagen externa.
        \item Determinar el umbral de binarización.
    \end{enumerate}

    % TEMA 4: Morfología
    \item \textbf{¿Qué operación morfológica utilizaría para rellenar pequeños agujeros oscuros dentro de un objeto blanco brillante?}
    \begin{enumerate}[label=\alph*)]
        \item Erosión.
        \item Apertura (Opening).
        \item Cierre (Closing).
        \item Esqueletonización.
    \end{enumerate}

    % TEMA 4: Clusterización
    \item \textbf{El algoritmo K-medias (K-means) es un método de segmentación:}
    \begin{enumerate}[label=\alph*)]
        \item Supervisado (requiere entrenamiento previo).
        \item No supervisado (agrupa píxeles en $k$ clases basándose en intensidad).
        \item Basado en detección de bordes mediante gradientes.
        \item Basado en atlas probabilísticos.
    \end{enumerate}

\end{enumerate}

% ---------------------------
% PARTE II: PROBLEMAS
% ---------------------------

\section*{II. Práctica (6 ptos)}
\textit{\footnotesize Resuelva los siguientes ejercicios justificando los pasos.}

% ---------------------------
% PROBLEMA 1
% ---------------------------

\textbf{P1. Filtros de Suavizado (2 ptos)}\\
Se tiene la siguiente sub-imagen original $I$ de $3 \times 3$, que presenta un píxel central con ruido notable respecto a sus vecinos:
\[
I = 
\begin{bmatrix}
10 & 10 & 10 \\
10 & \mathbf{90} & 10 \\
10 & 10 & 10
\end{bmatrix}
\]
Se desea aplicar un \textbf{filtro de media} (promedio) de tamaño $3 \times 3$. La máscara del filtro $H$ es:
\[
H = \frac{1}{9}
\begin{bmatrix}
1 & 1 & 1 \\
1 & 1 & 1 \\
1 & 1 & 1
\end{bmatrix}
\]

\begin{enumerate}[label=\alph*)]
    \item Calcule el nuevo valor del \textbf{píxel central} tras aplicar la convolución.
    \item Si en lugar del filtro de media usáramos un \textbf{filtro de mediana} $3 \times 3$, ¿cuál sería el valor del píxel central? Justifique cuál de los dos filtros elimina mejor este tipo de ruido impulsivo.
\end{enumerate}

\vspace{1mm}

% ---------------------------
% PROBLEMA 2
% ---------------------------

\textbf{P2. Transformación Geométrica (2 ptos)} \\
Consideramos un píxel en la coordenada $P(x, y) = (4, 10)$. 
\textbf{Nota:} Para todos los apartados, resuelva utilizando \textbf{coordenadas homogéneas} y exprese explícitamente las \textbf{matrices de transformación} utilizadas.

\begin{enumerate}
    \item[(a)] Defina la matriz de traslación $M_T$ para el vector $\vec{v} = (6, -5)$. Realice el producto matricial para obtener el punto transformado $P'$ en coordenadas homogéneas.
    
    \item[(b)] Partiendo de $P'$, defina la matriz de escalado $M_S$ para reducir la imagen a la mitad en ambos ejes ($s_x = 0.5, s_y = 0.5$). Calcule el nuevo punto $P''$ mediante multiplicación matricial.
    
    \item[(c)] Si quisiéramos rotar el punto original $P(4, 10)$ 90$^{\circ}$ en sentido \textbf{horario} alrededor del origen: escriba la matriz de rotación $M_R$ correspondiente y calcule la coordenada final.
\end{enumerate}

\vspace{1mm}

% ---------------------------
% PROBLEMA 3
% ---------------------------

\textbf{P3. Morfología Matemática: Erosión (2 ptos)}\\
Dada la siguiente imagen binaria $I$ de $5 \times 5$ (donde 1 es objeto y 0 es fondo):

\[
I = 
\begin{bmatrix}
0 & 0 & 0 & 0 & 0 \\
0 & 1 & 1 & 1 & 0 \\
0 & 1 & 1 & 1 & 0 \\
0 & 0 & 1 & 0 & 0 \\
0 & 0 & 0 & 0 & 0
\end{bmatrix}
\]

Se aplica una \textbf{Erosión} utilizando un elemento estructurante $B$ en forma de cruz ($3 \times 3$):
\[
B = 
\begin{bmatrix}
0 & 1 & 0 \\
1 & \mathbf{1} & 1 \\
0 & 1 & 0
\end{bmatrix}
\]
(El centro del elemento estructurante está marcado en negrita).

\begin{enumerate}[label=\alph*)]
    \item Calcule la imagen resultante tras la erosión.
    \item ¿Qué efecto visual ha tenido la erosión sobre el píxel aislado en la posición $(4,3)$ (fila 4, columna 3)?
\end{enumerate}

\end{multicols*}

\end{document}
